\section{Comparison theorem}\label{sec:comparison}

\subsection{The two fundamental properties}
\begin{paragr}\label{paragr:colimtr}
  Let $C$ be a $1$\nbd-category. Recall from \ref{paragr:pullback_resol} that, given an $\ook
  k$\nbd-category
  $X$ and an $\ook k$\nbd-functor $p \colon X \to C$, with $1 \leq k \leq \oo$, we defined a functor
    \[
    \begin{aligned}
      \tr{X}{-} \colon C &\to \opCat{k} \\
      c &\mapsto \tr{X}{c}.
    \end{aligned}
  \]
    This
    construction is functorial in $p\colon X \to C$ as it canonically defines a functor
  \begin{align*}
    \tr{\opCat{k}}{C} &\to \funopCat{k}{C}\\
    (X, X \overset{p}{\to} C) &\mapsto \left(c \mapsto \tr{X}{c}\right).
  \end{align*}
  Notice that for $(X, X \overset{p}{\to} C)$ in $\tr{\opCat{k}}{C}$,
  we have a cocone
  \[
    (\tr{X}{c} \to X)_{c \in \Ob{C}},
  \]
  hence
  a canonical morphism
  \[
    \ilim_{c \in C}\tr{X}{c} \to X,
  \]
  which is natural in the sense that it defines a natural
  transformation from the functor
  \begin{align*}
    \tr{\opCat{k}}{C} &\to \opCat{k}\\
    (X, X \overset{p}{\to} C) &\mapsto \ilim_{c \in C} \tr{X}{c},
  \end{align*}
  to the projection functor
    \begin{align*}
    \tr{\opCat{k}}{C} &\to \opCat{k}\\
    (X, X \overset{p}{\to} C) &\mapsto X.
    \end{align*}
    Finally, if we take $C$ a groupoid and $k=0$, then everything said
    above still makes sense.
\end{paragr}
\begin{lemma}\label{lemma:colim}
  The natural morphism of $\opCat{k}$
  \[
     \ilim_{c \in C}\tr{X}{c} \to X,
   \]
   is an isomorphism.
\end{lemma}
\begin{proof}
The case $k=\oo$ is \cite[Lemma 7.5]{guetta:homcat}. The general case
follows from the fact that $\opCat{k} \to \oCat$
  preserves colimits.
\end{proof}
\begin{proposition}\label{prop:hocolimtr}
  Let $C$ be a $1$\nbd-category. Then for any $1 \leq k \leq \oo$, the canonical morphism of
  $\Loc{\opCat{k}}$
  \[
    \hocolim_{c\in C }\tr{C}{c} \to C
    \]
  is an isomorphism. If $C$ is a groupoid, then this also works for
  $k=0$.
\end{proposition}
\begin{proof}
  Let $\frgp{P} \to C$ be a polygraphic resolution of $C$ in
  $\opCat{k}$ and \[\tr{\frgp P\!}{-} \colon C \to \opCat{k}\] obtained
  as in section \ref{subsec:polres}. We get a commutative diagram in
  $\Loc{\opCat{k}}$
  \[
    \begin{tikzcd}
      \displaystyle\hocolim_{c \in C}\tr{\frgp P\!}{c} \ar[d] \ar[r]& \displaystyle\ilim_{c \in
        C}\tr{\frgp P\!}{c} \ar[d] \ar[r] & \frgp P \ar[d] \\
      \displaystyle\hocolim_{c \in C}\tr{C}{c} \ar[r]&\displaystyle\ilim_{c \in C}\tr{C}{c}\ar[r] & C.
    \end{tikzcd}
  \]
  Note that both $\ilim_{c \in C}$ in the middle are the images by the
  localisation functor \[\opCat{k} \to \Loc{\opCat{k}}\]
  of the actual
  colimits taken in the category $\opCat{k}$. The right vertical
  arrow is an isomorphism by hypothesis, the horizontal arrows of the
  square on the right are isomorphisms by \cref{lemma:colim}. By
  stability of trivial fibrations under pullback, the natural transformation
  $\tr{\frgp P\!}{-} \to \tr{C}{-}$ is a pointwise weak equivalence. Since homotopy colimits
  send pointwise weak equivalences to isomorphisms (in the homotopy
  category), this proves that the left vertical arrow is an
  isomorphism. Finally, the top left
  horizontal arrow is an isomorphism
  because $\tr{\frgp P\!}{-}$ is cofibrant by \cref{thm:projresol} (or
  \cref{thm:projresol_grp} in the case $k=0$). By applying
  two-out-of-three twice, we get that all the remaining arrows of this diagram
  are isomorphisms too. 
\end{proof}
\begin{lemma}\label{lemma:liftoplax}
  Let $0 \leq k \leq \omega$ and consider a diagram as below in $\opCat{k}$, 
  \[
    \begin{tikzcd}[row sep=large]
      A' \ar[d,"p"]\ar[r,bend left,"u'"] \ar[r,bend right,"v'"']& B' \ar[d,"q"] \\
      A \ar[r,bend left,"u",""{name=A,below}] \ar[r,bend
      right,"v"',""{name=B,above}]& B,
      \ar[from=A,to=B,Rightarrow,"\alpha"]
    \end{tikzcd}
  \]
  where the front and back faces are commutative. If $q$ is a trivial
  fibration and $A'$ is cofibrant in the folk model structure on
  $\opCat{k}$, then there exists an oplax transformation $\alpha'
  \colon u' \Rightarrow v'$ making the whole diagram commute, which means
  that
  \[
    q\wsk \alpha' = \alpha \wsk p.
  \]
\end{lemma}
\begin{proof}
  The data of the given diagram amounts to the following commutative
  diagram in $\opCat{k}$
  \[
    \begin{tikzcd}
      \Sph{0} \grayk{k} A' \ar[r]\ar[d] \ar[rr,bend
      left,"{(u',v')}"]& \opifun{k}(\Dsk{1}) \grayk{k} A'
      \ar[d]& B' \ar[d,"q"] \\
      \Sph{0} \grayk{k} A \ar[r] & \opifun{k}(\Dsk{1}) \grayk{k} A \ar[r,"\alpha"]
      & B.
    \end{tikzcd}
  \]
  Since $A'$ is cofibrant and $\grayk{k}$ makes the folk model
  structure on $\opCat{k}$ a monoidal model structure \cite{ara_lucas:folmon}, the
  arrow $ \Sph{0}  \grayk{k} A' \to  \opifun{k}(\Dsk{1}) \grayk{k} A'$ is a
  cofibration. Since $q$ is a trivial fibration, there exists a
  morphism $\alpha' \colon \opifun{k}(\Dsk{1})  \grayk{k} A' \to B'$ making the
  whole diagram commute.
\end{proof}
\begin{proposition}\label{prop:hlgycontrcat}
  Let $C$ be a small category with a terminal object. Then, for any $1
  \leq k \leq \oo$ the
  canonical morphism of $\Derp{\Z}$
  \[
    \Hop{k}(C) \to \Z,
  \]
  induced by the canonical functor $C \to \Term$, is an
  isomorphism. If $C$ is a groupoid, then it is also an isomorphism
  for $k=0$.
\end{proposition}
\begin{proof}
  Let us denote by $c_0$ the terminal object of $C$. We then have a
  natural transformation
    \[
    \begin{tikzcd}
      C \ar[r,bend left,"\Unit{C}",""{name=A,below}]
      \ar[r,bend right,"\cst{c_0}"',""{name=B,above}] & C
      \ar[from=A,to=B,Rightarrow]
    \end{tikzcd}
  \]
  from the identity functor on $C$ to the constant functor with value
  $c_0$.  By \ref{paragr:gray}, this defines an oplax
  transformation. Now, let
  \[
    \frgp{P} \to C
  \]
  be a polygraphic resolution of $C$. Since trivial fibrations are
  surjective on objects, we may choose an object $x_0$ of $P$ whose image in
  $C$ is $c_0$. Thus, we have a diagram
  \[
    \begin{tikzcd}[row sep=huge]
      \frgp{P} \ar[d]\ar[r,bend left,"\Unit{\frgp{P}}"] \ar[r,bend
      right,"\cst{x_0}"']& \frgp{P} \ar[d] \\
      C \ar[r,bend left,"\Unit{C}",""{name=A,below}] \ar[r,bend
      right,"\cst{c_0}"',""{name=B,above}]& C,
      \ar[from=A,to=B,Rightarrow]
    \end{tikzcd}
  \]
  such that the front and back faces commute. Since
  $\frgp{P}$ is cofibrant, and since the vertical arrows are trivial
  fibrations, we can apply \cref{lemma:liftoplax}, which yields an
  oplax transformation from the identity \oo\nbd-functor on $P$ to
  the constant \oo\nbd-functor with value $x_0$.

  Finally, applying \cref{cor:abeltrans}, we get a chain homotopy
  between the identity morphism \[\abelk{k}(P) \to
    \abelk{k}(P)\] and the composition
  \[\abelk{k}(P) \to \Z \to
  \abelk{k}(P),\] which proves that  $\abelk{k}(P) \to
  \Z$ is a quasi-isomorphism.
\end{proof}
\subsection{The comparison theorem}
\begin{theorem}\label{thm:polhom}
  Let $C$ be a $1$\nbd-category. Then, for any $1
  \leq k \leq \oo$, there is a canonical isomorphism
  \[
    \Hop{k}(C) \to \Ho(C)
  \]
  natural in $C$. In particular, for $k=\oo$, we get an isomorphism
  $\Hopol(C) \to \Ho(C)$. Furthermore, for any $k\geq 1$, these
  isomorphisms make the following triangle commute
  \[
    \begin{tikzcd}[column sep=small]      
         \Hopol(C)\ar[dr]
         \ar[rr,"\homnat_C"]&&\Hop{k}(C)\ar[dl]\\
         &\Ho(C).&
    \end{tikzcd}
  \]
\end{theorem}

\begin{proof}
  Let us start with the first part of the theorem. Recall that if $F
  \colon \M \to \M'$ is a functor, and $C$ is a small category we
  denote by $\fun{C}{F} \colon \fun{C}{\M} \to \fun{C}{\M'}$ the
  functor induced by postcomposition. Consider the
  following zig-zag in $\Derp{\Z}$
  \[
    \begin{tikzcd}
      \hocolim_{C}p^*_C(\LL \abelk{k}(\Term))&&\\
    \hocolim_{C}\LL \fun{C}{\abelk{k}}(p_C^*\Term) \ar[u,"(a)"]&\ar[l,"(b)"'] \hocolim_{C}\LL\fun{C}{
     \abelk{k}}(\tr{C}{-})\ar[r,"(c)"]&  \LL
   \abelk{k}(\hocolim_{C}\tr{C}{-})\ar[d,"(d)"]\\
   && \LL\abelk{k}(C),
  \end{tikzcd}
  \]
  where
  \begin{itemize}
    \item $(a)$ and $(c)$ are the universal morphisms constructed in \ref{paragr:univmor},
    \item $(b)$ comes from the fact that $p_C^*\Term$ is the terminal
      object of $\funopCat{k}{C}$,
      \item $(d)$ is induced by the cocone $(\tr{C}{c}\to
  C)_{c \in \Ob(C)}$ (see \ref{paragr:colimtr}).
\end{itemize}
All these morphisms of $\Derp{\Z}$ are in fact isomorphisms:
\begin{itemize}
  \item $(a)$ by \eqref{eq:morphder} of \cref{prop:htpycocts} and the fact that $\abelk{k}$
    is a left Quillen functor (\cref{prop:abelquillen}),
      \item $(b)$ by \cref{prop:der2}, \cref{prop:hlgycontrcat} and the
    fact that for each object $c$ of $C$ the category $\tr{C}{c}$ has
    a terminal object,
  \item $(c)$ by \eqref{eq:cocontinuous} of \cref{prop:htpycocts} and the fact that $\abelk{k}$
    is a left Quillen functor,
  \item $(d)$ by \cref{prop:hocolimtr}.
  \end{itemize}
  Finally, since by
  \cref{paragr:abeldisk} and the fact that $\Dsk{0}=\Term$, we have
  \[
    \LL \abelk{k}(\Term) \cong \Z.
  \]
  Since, by definition, we have
  \[
    \Ho(C)=\hocolim_{C}p^*_C(\Z),
    \] this yields an isomorphism
  \[
    \Ho(C) \cong \Hop{k}(C),
  \]
  which is natural in $C$ by construction, as all the morphisms in the
  previous zig-zag are natural in $C$.

  For the second part, consider the following diagram in $\Derp{\Z}$
  \[
    \begin{tikzcd}
      \LL\abel(C)\ar[r] & \LL\abelk{k}(C)\\
       \LL\abel(\hocolim_{C}\tr{C}{-})\ar[r]\ar[u]&
       \LL\abelk{k}(\hocolim_{C}\tr{C}{-})\ar[u]\\
       \hocolim_{C}\fun{C}{\LL
     \abel}(\tr{C}{-})\ar[r] \ar[u]\ar[d]& \hocolim_{C}\fun{C}{\LL
     \abelk{k}}(\tr{C}{-})\ar[u]\ar[d]\\
   \hocolim_{C}\LL \fun{C}{\abel}(p_C^*\Term)\ar[r] \ar[d]& \hocolim_{C}\LL
   \fun{C}{\abelk{k}}(p_C^*\Term)\ar[d] \\
    \hocolim_{C}p^*_C(\LL \abel(\Term))\ar[r]& \hocolim_{C}p^*_C(\LL
    \abelk{k}(\Term)),
    \ar[from=1-1,to=2-2,phantom,description, "(A)"]
    \ar[from=2-1,to=3-2,phantom,description,"(B)"]
    \ar[from=3-1,to=4-2,phantom,description,"(C)"]
    \ar[from=4-1,to=5-2,phantom,description,"(D)"]
  \end{tikzcd}
\]
where the horizontal arrows are induced by the natural transformation
\eqref{transnat} from \cref{paragr:subtlety}
\[
    \homnat \colon \LL \abel\circ \oemb{k} \Rightarrow \LL \abelk{k}.
  \]
  Note that for simplicity we have not made the functor
  $\oemb{k} \colon \Loc{\opCat{k}}\to \Loc{\oCat}$ appear in the above
  diagram. This naturality implies that the squares $(A)$ and $(C)$ are
  commutative. Furthermore, the commutativity of the squares $(B)$ and
  $(D)$ follows from the naturality in $F$ of \eqref{eq:morphder} and
  \eqref{eq:cocontinuous} from \cref{prop:htpycocts}.

  To conclude, we only need to notice that the following triangle is
  commutative as an immediate computation shows
  \[
    \begin{tikzcd}
      \LL \abel(\Term) \ar[d,"\cong"']\ar[r] & \LL \abelk{k}(\Term)\ar[dl,"\cong"]\\
      \Z.&
    \end{tikzcd}
  \]
\end{proof}
\begin{corollary}
  For any $1$\nbd-category $C$ and any $1 \leq k \leq \oo$, the
  canonical natural morphism
  \[
   \homnat_C \colon \Hopol(C) \to \Hop{k}(C)
  \]
  is an isomorphism. If $C$ is a groupoid, this also holds for $k=0$.
\end{corollary}
\begin{remark}
  Adapting \emph{mutatis mutandis} the proof of
  \cref{thm:polhom}, one can prove more generally that, for a
  $1$-category $C$ and for any $1 \leq k \leq k' \leq \oo$, there is a canonical natural isomorphism
  \[
     \Hop{k'\!}(C)\cong  \Hop{k}(C).
  \]
\end{remark}

%%% Local Variables:
%%% mode: LaTeX
%%% TeX-master: "main"
%%% End:
